\chapter{Анализ предметной области и постановка задачи}
\label{ch:analysis}

\section{Характеристика предметной области}

Планирование пешеходного маршрута в городе включает три взаимосвязанных аспекта:
\begin{itemize}
  \item \textbf{Семантический} --- какие объекты городской среды (парки, набережные,
        велодорожки) пользователь считает привлекательными;
  \item \textbf{Комбинаторный} --- в каком порядке посещать выбранные точки между
        стартом и финишем;
  \item \textbf{Геометрический} --- как пройти по реально существующим тротуарам,
        дорожкам и аллеям, не пересекая непроходимые преграды.
\end{itemize}

Точки интереса (Points of Interest, POI) --- именованные объекты с координатами
и набором атрибутов. В проекте Fun-Walk каждая POI описывается типом
(\texttt{park}, \texttt{green\_zone}, \texttt{bike\_path}, \texttt{waterfront},
\texttt{square}), индексом качества воздуха AQI и уровнем шума (дБ). Демо-режим
использует 15~POI в центре Москвы: Парк Горького, Нескучный сад, Музеон, Крымская
набережная, Патриаршие пруды, фрагменты Садового кольца и др.

Пользователь задаёт \textbf{вектор предпочтений} --- шесть чисел от 0 до 10,
отражающих приоритет каждого критерия. Система должна вернуть маршрут с метриками:
длина (км), время (мин), зелёное покрытие (\%), балл качества (0--100), а также
список посещённых POI.

\section{Обзор существующих решений}

Рассмотрим классы аналогов (табл.~\ref{tab:analogs}).

\begin{table}[H]
\centering
\caption{Сравнительный анализ аналогов}
\label{tab:analogs}
\begin{tabular}{@{}p{3.2cm}p{4.5cm}p{4.5cm}p{3.5cm}@{}}
\toprule
\textbf{Решение} & \textbf{Сильные стороны} & \textbf{Ограничения} & \textbf{Открытый код} \\
\midrule
Google Maps & Точная маршрутизация, POI & Закрытый API, нет тонкой настройки «зелёности» & Нет \\
Komoot & Профили для пеших/велопрогулок & Коммерческая модель, облако & Нет \\
Walkonomics & Оценка «приятности» улиц & Устаревшие данные, нет API & Нет \\
OSRM + OSM & Свободная маршрутизация по OSM & Нет учёта предпочтений из коробки & Да \\
Fun-Walk (данная работа) & Явная матем. модель + UI & Демо-POI, in-memory хранение & Да \\
\bottomrule
\end{tabular}
\end{table}

Вывод: существует ниша для учебного и исследовательского приложения, где
математическая модель прозрачна, а исходный код доступен для модификации.

\section{Формальная постановка задачи}

Пусть заданы:
\begin{itemize}
  \item точка старта $S = (\varphi_S, \lambda_S)$ и точка финиша $E = (\varphi_E, \lambda_E)$;
  \item множество POI $\mathcal{P} = \{p_1, \ldots, p_n\}$;
  \item вектор предпочтений пользователя $\mathbf{w} \in \mathbb{R}^6_{\geq 0}$.
\end{itemize}

Требуется найти последовательность waypoints $W = (S, v_1, \ldots, v_k, E)$, где
$v_i \in \mathcal{P} \cup \emptyset$, и полилинию $\Gamma(W)$ по проходимым путям OSM,
максимизирующую интегральную оценку качества при ограничении на разумное
увеличение длины относительно кратчайшего пути $d_H(S, E)$.

Задача декомпозируется на две подзадачи:
\begin{enumerate}
  \item \textbf{Дискретная оптимизация} --- выбор подмножества POI и порядка обхода
        на абстрактном графе $G = (V, E)$ методом Dijkstra;
  \item \textbf{Непрерывная геометрия} --- построение $\Gamma$ через внешний
        маршрутизатор OSRM с профилем \texttt{foot}.
\end{enumerate}

Такое разделение снижает вычислительную сложность: $|V| \leq 10$, что делает
классический Dijkstra практичным без эвристик A*.

\section{Требования к программному комплексу}

На основании анализа сформулированы функциональные и нефункциональные требования.

\textbf{Функциональные:}
\begin{itemize}
  \item настройка шести приоритетов через UI;
  \item выбор старта и финиша кликом на карте;
  \item построение маршрута по REST API;
  \item отображение polyline, маркеров POI и метрик;
  \item сохранение, просмотр и удаление маршрутов в рамках сессии сервера.
\end{itemize}

\textbf{Нефункциональные:}
\begin{itemize}
  \item кросс-платформенность (Windows, Linux);
  \item время ответа API $\leq 15$~с (ограничение OSRM timeout);
  \item модульность backend (NestJS) и типизация (TypeScript);
  \item graceful degradation при недоступности OSRM.
\end{itemize}

\section{Выбор технологического стека}

Backend реализован на NestJS~10 --- фреймворке для Node.js с модульной архитектурой,
внедрением зависимостей и декораторами. Frontend --- SPA на React~18 и Vite~6.
Картография: Leaflet + тайлы OpenStreetMap. Маршрутизация: публичный OSRM-сервер
(\texttt{router.project-osrm.org}). Выбор TypeScript обеспечивает единую типовую
модель на клиенте и сервере (\texttt{LatLng}, \texttt{RoutePreferences}).

Обоснование: стек широко используется в индустрии, имеет обширную документацию и
позволяет сосредоточиться на алгоритмической части, а не на инфраструктурных
деталях.

\section{Выводы по главе}

Предметная область охватывает многокритериальную оценку POI, комбинаторный выбор
маршрута и геометрическую привязку к дорожной сети. Постановка задачи сведена к
двухэтапному алгоритму: Dijkstra на малом графе + OSRM. Сформулированы требования
и обоснован выбор стека NestJS/React/TypeScript.
